% GENERATED FILE. Edit canonical lessons, then run npm run generate:books.
% canonical-source-hash: fnv1a64:7a2f4ed232969938
% canonical-lessons: AR-C283-yaktubu, AR-C283-katabtu, AR-C283-kana, AR-C283-kana-yaktubu, AR-R283-the-past-person-by-person

\chapter{The Past, Person By Person}
\label{ch:ar-the-past-person-by-person}
\hlchaptermodality{283}

\begin{chapteropening}
\textbf{By the end of this chapter:} \emph{I can say what I do now and what I did, for every person, and paint the scene of a story with kāna and the present.}

Complete AR-R283-the-past-person-by-person: a scene with kāna and the present, and the event that cut into it, joined with ḥīnamā.
\end{chapteropening}

\section[\texorpdfstring{\ar{يكتب} (yaktubu)}{يكتب (yaktubu)}]{\ar{يكتب} (yaktubu) — the present, for everyone}

\label{lesson:AR-C283-yaktubu}



\begin{quote}
Say \textquotedblleft{}I write\textquotedblright{}. (\emph{Aktub.}) Then say the Arabic for a sentence and a barber. (\emph{Jumla}, \emph{ḥallāq}.)

Chapter thirty-seven gave you the present for one person only: \textbf{\ar{أ}} at the front means \emph{I}. A story needs \emph{he}, \emph{she} and \emph{we} as well, and they come from the same front door.
\end{quote}

\begin{grammarlens}[title={Grammar Lens: one stem, four prefixes}]
Keep the stem \emph{-ktub-} and change the letter in front of it:

\begin{quote}
\textbf{\ar{أكتب}} \emph{aktubu}, I write · \textbf{\ar{تكتب}} \emph{taktubu}, you write · \textbf{\ar{يكتب}} \emph{yaktubu}, he writes · \textbf{\ar{تكتب}} \emph{taktubu}, she writes · \textbf{\ar{نكتب}} \emph{naktubu}, we write
\end{quote}

\emph{Ta-} does two jobs: \textquotedblleft{}you\textquotedblright{} said to a man, and \textquotedblleft{}she\textquotedblright{}. The sentence around it tells them apart. Said to a woman, \textquotedblleft{}you write\textquotedblright{} adds \emph{-īna} at the back, \emph{taktubīna}; \textquotedblleft{}they write\textquotedblright{} adds \emph{-ūna}, \emph{yaktubūna}, and so does \textquotedblleft{}you write\textquotedblright{} said to several people, \emph{taktubūna}.

The \emph{-u} at the end is the present's own vowel. At the end of a sentence it is usually dropped, which is why chapter thirty-seven wrote \emph{aktub}; inside a sentence you will hear it.
\end{grammarlens}

\begin{grammarlens}[title={Grammar Lens: the same door on other verbs}]
\begin{quote}
\textbf{\ar{هي} \ar{تكتب} \ar{رسالة}.} — \emph{hiya taktubu risāla} — She is writing a letter.
\end{quote}

\begin{quote}
\textbf{\ar{نحن} \ar{نذهب} \ar{إلى} \ar{المدرسة}.} — \emph{naḥnu nadhhabu ilā al-madrasa} — We go to school.
\end{quote}

\emph{Dhahaba} gives \emph{nadhhabu}; \emph{qaraʾa} gives \emph{yaqraʾu}, he reads. One present covers \textquotedblleft{}I write\textquotedblright{} and \textquotedblleft{}I am writing\textquotedblright{}; English splits what Arabic keeps as one. \emph{Anā} is \textquotedblleft{}I\textquotedblright{}, \emph{huwa} \textquotedblleft{}he\textquotedblright{}, \emph{hiya} \textquotedblleft{}she\textquotedblright{}, \emph{naḥnu} \textquotedblleft{}we\textquotedblright{}, and they are often left out, because the prefix already says who.
\end{grammarlens}

\subsection*{Your turn}
Say these aloud:

\begin{itemize}
\raggedright
  \item \emph{aktubu, taktubu, yaktubu, naktubu}
  \item \textquotedblleft{}She is writing a letter\textquotedblright{} — \emph{hiya taktubu risāla}
  \item \textquotedblleft{}He reads a book\textquotedblright{} — \emph{yaqraʾu kitāban}
  \item \textquotedblleft{}We go to school\textquotedblright{} — \emph{nadhhabu ilā al-madrasa}
  \item \textquotedblleft{}The actor reads\textquotedblright{}, with the word from chapter two hundred and eighty-one — \emph{al-mumaththil yaqraʾu}
\end{itemize}

\begin{tcolorbox}[breakable,colback=teal!4,colframe=teal!35!black,title={Before you move on}]
Which letter at the front means \emph{we}? (\textbf{\ar{ن}}, \emph{na-}.) And \emph{he}? (\textbf{\ar{ي}}, \emph{ya-}.) Why can \textbf{\ar{تكتب}} mean two people? (\emph{Ta-} is \textquotedblleft{}you\textquotedblright{} and \textquotedblleft{}she\textquotedblright{}: \textbf{the sentence decides}.) Where is the \emph{-u} dropped? (\textbf{At the end of a sentence.})
\end{tcolorbox}
\section[\texorpdfstring{\ar{كتبت} (katabtu)}{كتبت (katabtu)}]{\ar{كتبت} (katabtu) — the past, person by person}

\label{lesson:AR-C283-katabtu}



\begin{quote}
Say \textquotedblleft{}we write\textquotedblright{}, then \textquotedblleft{}he writes\textquotedblright{}. (\emph{Naktubu}, \emph{yaktubu}.)

You have said \emph{qaḍaytu al-ʿuṭla fī al-qarya}, \textquotedblleft{}I spent the holiday in the village\textquotedblright{}. That \emph{-tu} at the back is \textquotedblleft{}I\textquotedblright{} in the past. Here is the rest of the set.
\end{quote}

\begin{grammarlens}[title={Grammar Lens: the endings at the back}]
\emph{Kataba}, he wrote, is the bare past. Every other person is an ending on the same stem, \emph{katab-}:

\begin{quote}
\textbf{\ar{كتبتُ}} \emph{katabtu}, I wrote · \textbf{\ar{كتبتَ}} \emph{katabta}, you wrote (to a man) · \textbf{\ar{كتبتِ}} \emph{katabti}, you wrote (to a woman) · \textbf{\ar{كتبَ}} \emph{kataba}, he wrote · \textbf{\ar{كتبتْ}} \emph{katabat}, she wrote · \textbf{\ar{كتبنا}} \emph{katabnā}, we wrote · \textbf{\ar{كتبوا}} \emph{katabū}, they wrote · \textbf{\ar{كتبتم}} \emph{katabtum}, you wrote (to several)
\end{quote}

The present put the person at the front; the past puts it at the back.

Without its vowel marks, \textbf{\ar{كتبت}} spells four of these at once: \emph{katabtu}, \emph{katabta}, \emph{katabti} and \emph{katabat}. That is why they carry marks here, and why, on an unmarked page, the rest of the sentence tells you who wrote.
\end{grammarlens}

\begin{grammarlens}[title={Grammar Lens: other verbs, and \textquotedblleft{}did not\textquotedblright{}}]
\begin{quote}
\textbf{\ar{أمس} \ar{كتبت} \ar{رسالة} \ar{إلى} \ar{جدتي}.} — \emph{ams katabtu risālatan ilā jaddatī} — Yesterday I wrote a letter to my grandmother.
\end{quote}

\begin{quote}
\textbf{\ar{ذهبنا} \ar{إلى} \ar{السوق}.} — \emph{dhahabnā ilā as-sūq} — We went to the market.
\end{quote}

\emph{Qaraʾa} gives \emph{qaraʾtu}, I read; \emph{dhahaba} gives \emph{dhahabat}, she went. The \textbf{\ar{ا}} at the end of \textbf{\ar{كتبوا}} is written and not said.

\textbf{\ar{ما}} in front makes it \textquotedblleft{}did not\textquotedblright{}: \emph{mā katabtu}, I didn't write; \emph{mā dhahabnā}, we didn't go.
\end{grammarlens}

\subsection*{Your turn}
Say these aloud:

\begin{itemize}
\raggedright
  \item \emph{katabtu, katabta, katabti, kataba, katabat, katabnā, katabū}
  \item \textquotedblleft{}Yesterday I wrote a letter\textquotedblright{} — \emph{ams katabtu risālatan}
  \item \textquotedblleft{}We went to the market\textquotedblright{}, then \textquotedblleft{}we didn't go\textquotedblright{} — \emph{dhahabnā ilā as-sūq}, \emph{mā dhahabnā}
  \item \textquotedblleft{}They came back to the house\textquotedblright{} — \emph{rajaʿū ilā al-manzil}
  \item \textquotedblleft{}The engineer read the letter\textquotedblright{} — \emph{al-muhandis qaraʾa ar-risāla}
\end{itemize}

\begin{tcolorbox}[breakable,colback=teal!4,colframe=teal!35!black,title={Before you move on}]
Where does the past mark the person? (\textbf{At the back.}) \textquotedblleft{}She wrote\textquotedblright{}? (\emph{\textbf{Katabat.}}) \textquotedblleft{}We wrote\textquotedblright{}? (\emph{\textbf{Katabnā.}}) How many persons can an unmarked \textbf{\ar{كتبت}} be? (\textbf{Four}: \emph{katabtu, katabta, katabti, katabat}.) \textquotedblleft{}I didn't go\textquotedblright{}? (\emph{\textbf{Mā dhahabtu.}})
\end{tcolorbox}
\section[\texorpdfstring{\ar{كان} (kāna)}{كان (kāna)}]{\ar{كان} (kāna) — was}

\label{lesson:AR-C283-kana}



\begin{quote}
Say \textquotedblleft{}I wrote\textquotedblright{} and \textquotedblleft{}she wrote\textquotedblright{}. (\emph{Katabtu}, \emph{katabat}.)

Now say \textquotedblleft{}I am at the office\textquotedblright{}. (\emph{Anā fī al-maktab}.) There is no \textquotedblleft{}am\textquotedblright{} in it: in the present, Arabic sets the two halves side by side. In the past it cannot, and it needs a verb.
\end{quote}

\subsection*{You'll want to know: \ar{كان}}
\textbf{\ar{كان}} — \emph{kāna} — \textquotedblleft{}he was, it was\textquotedblright{}: the past of \textquotedblleft{}to be\textquotedblright{}.

\begin{quote}
\textbf{\ar{كنت} \ar{في} \ar{المكتب} \ar{أمس}.} — \emph{kuntu fī al-maktab ams} — I was at the office yesterday.
\end{quote}

\begin{quote}
\textbf{\ar{كانت} \ar{جدتي} \ar{في} \ar{القرية}.} — \emph{kānat jaddatī fī al-qarya} — My grandmother was in the village.
\end{quote}

It takes the endings of the lesson before, with one change. Where the ending starts with a consonant, the long \emph{ā} shrinks to \emph{u}: \emph{kuntu}, I was; \emph{kunta}, \emph{kunti}, you were; \emph{kunnā}, we were. Where it does not, the \emph{ā} stays: \emph{kāna}, \emph{kānat}, \emph{kānū}, they were.

\begin{cousinweb}
\textbf{\ar{كان}} (\emph{kāna}) is built on the native Arabic root \emph{k-w-n}, \textquotedblleft{}to be, to come into being\textquotedblright{} (Lane). Lane lists \emph{makāna}, \textquotedblleft{}a place of being\textquotedblright{}, among the words of the same root. The \emph{w} of the root is the reason for the two shapes: it surfaces as long \emph{ā} in \emph{kāna} and as short \emph{u} in \emph{kuntu}.
\end{cousinweb}

\subsection*{Your turn}
Say these aloud:

\begin{itemize}
\raggedright
  \item \emph{kāna, kānat} — then \emph{kuntu, kunnā}
  \item \textquotedblleft{}I was at the office yesterday\textquotedblright{} — \emph{kuntu fī al-maktab ams}
  \item \textquotedblleft{}We were in the restaurant\textquotedblright{} — \emph{kunnā fī al-maṭʿam}
  \item \textquotedblleft{}My grandmother was in the village\textquotedblright{} — \emph{kānat jaddatī fī al-qarya}
\end{itemize}

\begin{tcolorbox}[breakable,colback=teal!4,colframe=teal!35!black,title={Before you move on}]
What does \ar{كان} mean? (\textquotedblleft{}He was; it was\textquotedblright{}.) \textquotedblleft{}I was\textquotedblright{}? (\emph{\textbf{Kuntu}}, with a short \emph{u}.) \textquotedblleft{}She was\textquotedblright{}? (\emph{\textbf{Kānat}}, with a long \emph{ā}.) Why does the present not need it? (\textbf{Arabic sets the two halves side by side.})
\end{tcolorbox}
\section[\texorpdfstring{\ar{كان} \ar{يكتب} (kāna yaktubu)}{كان يكتب (kāna yaktubu)}]{\ar{كان} \ar{يكتب} (kāna yaktubu) — the past that is still running}

\label{lesson:AR-C283-kana-yaktubu}



\begin{quote}
Say \textquotedblleft{}I was at the office\textquotedblright{}, then \textquotedblleft{}I wrote\textquotedblright{}. (\emph{Kuntu fī al-maktab}, \emph{katabtu}.)

\emph{Katabtu} has an edge: the writing happened, and it is over. A story also needs a past with no edge on it.
\end{quote}

\begin{grammarlens}[title={Grammar Lens: kāna in front of the present}]
Put \emph{kāna} in its past form, then the present of the verb you mean, both for the same person:

\begin{quote}
\emph{kuntu} + \emph{aktubu} $\to$ \textbf{\ar{كنت} \ar{أكتب}} — \emph{kuntu aktubu} — I was writing; I used to write
\end{quote}

So \emph{kānat taqraʾu}, she was reading; \emph{kunnā nadhhabu}, we used to go. It answers two English shapes. \textquotedblleft{}Was ‑ing\textquotedblright{}: the action was under way. \textquotedblleft{}Used to\textquotedblright{}: it happened again and again.

\begin{quote}
\textbf{\ar{كنت} \ar{أمشي} \ar{إلى} \ar{المدرسة} \ar{كل} \ar{يوم}.} — \emph{kuntu amshī ilā al-madrasa kulla yawm} — I used to walk to school every day.
\end{quote}

\emph{Mā} goes in front of \emph{kāna}: \emph{mā kuntu aʿrifu}, I didn't know.
\end{grammarlens}

\begin{grammarlens}[title={Grammar Lens: the scene, and the event}]
A told story splits its verbs in two. \emph{Kāna} with the present paints the \textbf{scene}: what was going on. The plain past drops the \textbf{events} into it: what happened.

\begin{quote}
\textbf{\ar{كنت} \ar{أقرأ} \ar{كتابا} \ar{حينما} \ar{اتصل} \ar{صديقي}.} — \emph{kuntu aqraʾu kitāban ḥīnamā ittaṣala ṣadīqī} — I was reading a book when my friend phoned.
\end{quote}

The reading is scene, still running when the call cuts into it. The call is an event, over in a moment, and \emph{ḥīnamā} is the joint between them.
\end{grammarlens}

\subsection*{Your turn}
Say these aloud:

\begin{itemize}
\raggedright
  \item \emph{kuntu aktubu, kānat taktubu, kunnā naktubu}
  \item \textquotedblleft{}I used to walk to school every day\textquotedblright{} — \emph{kuntu amshī ilā al-madrasa kulla yawm}
  \item \emph{kuntu aqraʾu kitāban ḥīnamā ittaṣala ṣadīqī}
  \item \textquotedblleft{}I didn't know\textquotedblright{} — \emph{mā kuntu aʿrifu}
  \item \textquotedblleft{}She was reading a page\textquotedblright{} — \emph{kānat taqraʾu ṣafḥa}
\end{itemize}

\begin{tcolorbox}[breakable,colback=teal!4,colframe=teal!35!black,title={Before you move on}]
How do you build \textquotedblleft{}I was writing\textquotedblright{}? (\emph{Kuntu} + \emph{aktubu}: \emph{\textbf{kuntu aktubu}}.) In the phone sentence, which verb is the scene? (\emph{\textbf{Kuntu aqraʾu}}, still running.) Which is the event? (\emph{\textbf{Ittaṣala}}, over in a moment.) Do both verbs agree with the person? (\textbf{Yes}: \emph{kunnā naktubu}.)
\end{tcolorbox}
\section[(review)]{Review — the scene, and what cut into it}

\label{lesson:AR-R283-the-past-person-by-person}



\begin{quote}
Say \textquotedblleft{}he writes\textquotedblright{} for today, then \textquotedblleft{}he wrote\textquotedblright{} for yesterday. (\emph{Yaktubu}, \emph{kataba}.) Then the Arabic for a calculator and a pencil case. (\emph{Ḥāsiba}, \emph{miqlama}.)
\end{quote}

\begin{grammarlens}[title={What you've built}]
Four lessons, one stem, and three shapes of it. The present marks the person at the front: \emph{aktubu, taktubu, yaktubu, naktubu}. The past marks it at the back: \emph{katabtu, katabat, katabnā, katabū}. And \emph{kāna} in front of the present gives the past that is still running: \emph{kuntu aktubu}.

\emph{Al-māḍī} sets the time:

\begin{quote}
\textbf{\ar{كنت} \ar{أعمل} \ar{في} \ar{المكتب} \ar{الأسبوع} \ar{الماضي}.} — \emph{kuntu aʿmalu fī al-maktab al-usbūʿa al-māḍī} — Last week I was working at the office.
\end{quote}

And \emph{ḥīnamā} joins the scene to the event that cut into it:

\begin{quote}
\textbf{\ar{كانت} \ar{تكتب} \ar{قصة} \ar{حينما} \ar{جاء} \ar{أخوها}.} — \emph{kānat taktubu qiṣṣa ḥīnamā jāʾa akhūhā} — She was writing a story when her brother came.
\end{quote}
\end{grammarlens}

\subsection*{Your turn}
Say these aloud:

\begin{itemize}
\raggedright
  \item a scene in two — \emph{kuntu fī al-maktab. kuntu aktubu.}
  \item \textquotedblleft{}Last week I was working at the office\textquotedblright{} — \emph{kuntu aʿmalu fī al-maktab al-usbūʿa al-māḍī}
  \item \emph{kānat taktubu qiṣṣa ḥīnamā jāʾa akhūhā}
  \item \textquotedblleft{}The pilot was waiting\textquotedblright{}, then \textquotedblleft{}the presenter came\textquotedblright{} — \emph{kāna aṭ-ṭayyār yantaẓiru}, \emph{jāʾa al-mudhīʿ}
  \item one of your own: \emph{kuntu} with a present, then \emph{ḥīnamā} and one thing that happened
\end{itemize}

\begin{tcolorbox}[breakable,colback=teal!4,colframe=teal!35!black,title={Before you move on}]
After \emph{ḥīnamā} in a story, which past? (The plain one, the \textbf{event}: \emph{jāʾa}.) Which past paints the scene? (\emph{\textbf{Kāna}} with the present.) \textquotedblleft{}We were\textquotedblright{}? (\emph{\textbf{Kunnā.}}) \textquotedblleft{}Last week\textquotedblright{}? (\emph{\textbf{Al-usbūʿ al-māḍī.}})
\end{tcolorbox}
